% !TeX root = ../main.tex
% Contenuto della Lezione 8 importato dalla trascrizione già revisionata.
% Convenzioni e preambolo sono ora centralizzati; i tre TikZ sono conservati.
% L'introduzione dalla Lezione 7 precede le proprietà della Lezione 8.
\chapter{Topologia della retta}
\label{cap:topologia}
\fonti{Lezioni 7 e 8}

% !TeX root = ../main.tex
% Fonte: Lezione 7.pdf, pp. 19--21.
\section{Insiemi aperti e insiemi chiusi}

\begin{definizione}[Insieme aperto]
Un insieme $A\subseteq\R$ si dice \emph{aperto} se è un'unione di intervalli aperti: esiste una famiglia $(I_\alpha)_{\alpha\in J}$ di intervalli aperti tale che
\[
A=\bigcup_{\alpha\in J}I_\alpha.
\]
\end{definizione}

\textbf{Famiglie ammesse.} La famiglia può avere un numero finito o infinito di elementi; si ammette anche la famiglia vuota, la cui unione è $\varnothing$.

\textbf{Esempi di aperti.}
\begin{itemize}
  \item $\R=(-\infty,+\infty)$ è aperto;
  \item $\varnothing$ è aperto, come unione della famiglia vuota;
  \item $\R\setminus\Z=\bigcup_{n\in\Z}(n,n+1)$ è aperto;
  \item $\R^*=\R\setminus\{0\}=(-\infty,0)\cup(0,+\infty)$ è aperto.
\end{itemize}
L'insieme di tutti gli aperti si chiama \emph{topologia usuale} di $\R$.

\begin{definizione}[Insieme chiuso]
Un insieme $B\subseteq\R$ si dice \emph{chiuso} se il complemento $B^c=\R\setminus B$ è aperto.
\end{definizione}

\textbf{Esempi di chiusi.} Per verificare che un insieme è chiuso basta quindi descrivere il complemento come unione di intervalli aperti:
\begin{itemize}
  \item $\R$ è chiuso, perché $\R^c=\varnothing$ è aperto;
  \item $\varnothing$ è chiuso, perché $\varnothing^c=\R$ è aperto;
  \item per $a\le b$, $[a,b]$ è chiuso, perché $[a,b]^c=(-\infty,a)\cup(b,+\infty)$;
  \item $\{0\}$ è chiuso, perché il suo complemento è $(-\infty,0)\cup(0,+\infty)$;
  \item $[a,+\infty)$ è chiuso, perché ha complemento $(-\infty,a)$.
\end{itemize}
\textbf{Rapporto fra le due proprietà.} Aperto e chiuso non sono proprietà mutuamente esclusive: $\varnothing$ e $\R$ le possiedono entrambe. Esistono anche insiemi che non possiedono nessuna delle due, come vedremo per $[0,1)$.

\section{Il criterio locale di apertura}

\begin{proposizione}[Caratterizzazione degli aperti]
\label{prop:aperti-criterio}
Un insieme $A\subseteq\R$ è aperto se e solo se
\[
\forall x\in A\ \exists\varepsilon>0:
\quad (x-\varepsilon,x+\varepsilon)\subseteq A.
\]
\end{proposizione}
\begin{proof}
\passomath{\Rightarrow} Se $A$ è un'unione di intervalli aperti e $x\in A$, allora $x$ appartiene a uno di tali intervalli, diciamo $I$. Poiché $x$ non è un estremo di $I$, si può scegliere $\varepsilon>0$ abbastanza piccolo affinché $(x-\varepsilon,x+\varepsilon)\subseteq I\subseteq A$. Per $I=(a,b)$ con estremi finiti basta, per esempio, $\varepsilon=\min\{x-a,b-x\}/2$; per una semiretta basta considerare la distanza dall'unico estremo finito, e per $I=\R$ ogni raggio va bene.

\passomath{\Leftarrow} Per ogni $x\in A$ scegliamo un intervallo $I_x=(x-\varepsilon_x,x+\varepsilon_x)\subseteq A$ come nell'ipotesi. Allora
\[
A=\bigcup_{x\in A}I_x:
\]
\passomath{A\subseteq\bigcup_{x\in A}I_x} Ogni $x\in A$ appartiene al proprio intervallo $I_x$.

\passomath{\bigcup_{x\in A}I_x\subseteq A} Tutti gli intervalli $I_x$ sono contenuti in $A$. È dunque una rappresentazione di $A$ come unione di intervalli aperti. Se $A=\varnothing$, entrambi i criteri valgono con la convenzione della famiglia vuota.
\end{proof}

La proprietà è locale: intorno a ogni punto di $A$ si deve poter ritagliare un piccolo intervallo che resti interamente in $A$. Il raggio può dipendere dal punto; non serve un unico raggio valido per tutti.

\begin{esempio}
\textbf{Non apertura.} L'insieme $A=[0,1)$ non è aperto: $0\in A$, ma ogni intervallo $(-\varepsilon,\varepsilon)$ contiene punti negativi, esterni ad $A$.

\textbf{Non chiusura.} Non è neppure chiuso, perché
\[
A^c=(-\infty,0)\cup[1,+\infty)
\]
non è aperto: contiene $1$, ma nessun intervallo aperto centrato in $1$ è contenuto nel complemento. Basta infatti scegliere, per ogni $\varepsilon>0$, il punto $1-\min\{\varepsilon/2,1/2\}\in A\cap(1-\varepsilon,1+\varepsilon)$.
\end{esempio}


\section{Proprietà degli insiemi aperti}
% Fonte: Lezione 8, pp.~2--3

Ricordiamo il criterio fondamentale: un insieme $A\subseteq\R$ è \emph{aperto}
se, per ogni $x\in A$, esiste $\varepsilon>0$ tale che
$(x-\varepsilon,x+\varepsilon)\subseteq A$. In altre parole, ogni punto di $A$
può essere circondato da un piccolo intervallo interamente contenuto in $A$.

\begin{proposizione}[Operazioni con gli aperti]
Valgono le seguenti proprietà.
\begin{enumerate}
  \item $\varnothing$ e $\R$ sono aperti.
  \item L'intersezione di un numero \emph{finito} di aperti è aperta.
  \item L'unione di una famiglia \emph{arbitraria} di aperti è aperta.
\end{enumerate}
\end{proposizione}

\begin{proof}
Per l'insieme vuoto non ci sono punti sui quali verificare il criterio;
per $\R$, qualunque intervallo centrato in un suo punto resta in $\R$.

Siano $A_1,\ldots,A_m$ aperti e sia $x\in\bigcap_{j=1}^{m}A_j$.
Per ogni $j$ esiste $\varepsilon_j>0$ per cui
$(x-\varepsilon_j,x+\varepsilon_j)\subseteq A_j$.
Poiché gli insiemi sono in numero finito, il minimo
$\varepsilon=\min\{\varepsilon_1,\ldots,\varepsilon_m\}$ è ancora positivo.
L'intervallo $(x-\varepsilon,x+\varepsilon)$ è quindi contenuto in tutti
gli $A_j$, e dunque nella loro intersezione.

Se invece $x$ appartiene a un'unione $\bigcup_{\alpha\in I}A_\alpha$ di aperti,
allora appartiene ad almeno uno di essi, diciamo $A_{\alpha_0}$. Un intervallo
aperto intorno a $x$ contenuto in $A_{\alpha_0}$ è automaticamente contenuto
nell'intera unione. Non serve che la famiglia sia finita.
\end{proof}

L'ipotesi di finitezza nell'intersezione è essenziale. Per esempio,
\[
  \bigcap_{n\in\N^*}\left(-\frac1n,\frac1n\right)=\{0\}.
\]
Il numero $0$ appartiene a tutti gli intervalli. Se $x\ne0$, scegliendo
$n>1/|x|$ si ottiene $1/n<|x|$, perciò $x$ non appartiene all'$n$-esimo
intervallo. Tuttavia $\{0\}$ non è aperto: nessun intervallo aperto centrato
in $0$ può essere contenuto nel solo punto $0$. Un'intersezione infinita
di aperti può quindi non essere aperta.

\Needspace{12\baselineskip}
\section{Proprietà degli insiemi chiusi}
% Fonte: Lezione 8, p.~4

Un insieme $C\subseteq\R$ è \emph{chiuso} quando il suo complemento $C^c$
è aperto. Applicando le leggi di De Morgan alle proprietà appena dimostrate
si ottiene il risultato duale.

\begin{proposizione}[Operazioni con i chiusi]
\begin{enumerate}
  \item $\varnothing$ e $\R$ sono chiusi.
  \item L'intersezione di una famiglia arbitraria di chiusi è chiusa.
  \item L'unione di un numero finito di chiusi è chiusa.
\end{enumerate}
\end{proposizione}

\begin{proof}
I complementi di $\varnothing$ e $\R$ sono rispettivamente $\R$ e
$\varnothing$, che sono aperti. Se $C_\alpha$ sono chiusi, allora
\[
  \left(\bigcap_{\alpha\in I}C_\alpha\right)^c
  =\bigcup_{\alpha\in I}C_\alpha^c
\]
è aperto, perché è un'unione arbitraria di aperti. Analogamente, per
$C_1,\ldots,C_m$ chiusi,
\[
  \left(\bigcup_{j=1}^{m}C_j\right)^c
  =\bigcap_{j=1}^{m}C_j^c
\]
è aperto, perché l'intersezione è finita. In entrambi i casi l'insieme
di partenza è dunque chiuso.
\end{proof}

\section{Esempi nella retta reale}
% Fonte: Lezione 8, pp.~5--9

\subsection{I numeri naturali e gli interi}

L'insieme $\N$ è chiuso. Un modo diretto per vederlo è scrivere il suo
complemento come unione di intervalli aperti:
\[
  \R\setminus\N=(-\infty,0)\ \cup\
  \bigcup_{n\in\N}(n,n+1).
\]
Ogni reale non naturale è negativo oppure si trova strettamente tra due
naturali consecutivi. Il membro destro è aperto, essendo un'unione di aperti;
per definizione, quindi, $\N$ è chiuso.

Lo stesso ragionamento vale per gli interi:
\[
  \R\setminus\Z=\bigcup_{n\in\Z}(n,n+1).
\]
Fra due interi consecutivi non si trovano altri interi, e ogni reale non
intero cade in uno di questi intervalli. Ne segue che $\Z$ è chiuso.
Né $\N$ né $\Z$ sono aperti: ciascuno dei loro punti è isolato rispetto
all'insieme, cioè ogni piccolo intervallo centrato in esso contiene anche
reali che non gli appartengono.

\subsection{I numeri razionali}

$\Q$ non è aperto. La ragione è che in ogni intervallo aperto non vuoto
si trova almeno un numero irrazionale: dunque non esiste alcun intervallo
aperto centrato in un razionale e interamente contenuto in $\Q$.
Questa proprietà si chiama \emph{densità degli irrazionali in $\R$}.

Per chiarire il passaggio, siano $a<b$. Dalla densità di $\Q$ in $\R$
possiamo scegliere un razionale $q\ne0$ tale che
$a/\sqrt2<q<b/\sqrt2$. Allora $a<q\sqrt2<b$ e $q\sqrt2$ è irrazionale:
se fosse razionale, dividendo per il razionale non nullo $q$ si otterrebbe
che $\sqrt2$ è razionale. Abbiamo così trovato un irrazionale in $(a,b)$.

Neppure $\Q$ è chiuso. Infatti, se lo fosse, il complemento
$\R\setminus\Q$ sarebbe aperto. Ma, per la densità dei razionali,
ogni intervallo aperto centrato in un irrazionale contiene razionali;
quindi $\R\setminus\Q$ non è aperto. In sintesi, $\Q$ non è né aperto
né chiuso. La densità di entrambi gli insiemi spiega le due conclusioni.

\subsection{La successione \texorpdfstring{$1/n$}{1/n} vista come insieme}

Consideriamo
\[
  X=\left\{\frac1n:n\in\N^*\right\}
   =\left\{1,\frac12,\frac13,\ldots\right\}.
\]
$X$ non è aperto: per esempio, ogni intervallo aperto centrato in $1$
contiene numeri diversi da tutti gli elementi di $X$.

$X$ non è neppure chiuso. Il punto $0$ non appartiene a $X$, ma è
arbitrariamente vicino ai suoi elementi: per ogni $\varepsilon>0$ esiste
$n\in\N^*$ con $1/n<\varepsilon$. Il complemento di $X$ contiene dunque
$0$, ma non contiene nessun intervallo aperto centrato in $0$;
di conseguenza non è aperto. Esplicitamente,
\[
  X^c=(-\infty,0]\ \cup\
  \bigcup_{n\in\N^*}\left(\frac1{n+1},\frac1n\right)
  \ \cup\ (1,+\infty).
\]
L'estremo $0$, incluso nel primo pezzo, è proprio ciò che impedisce a
$X^c$ di essere aperto.

Se aggiungiamo il punto mancante, otteniamo invece l'insieme chiuso
\[
  Y=X\cup\{0\}.
\]
Infatti
\[
  Y^c=(-\infty,0)\ \cup\
  \bigcup_{n\in\N^*}\left(\frac1{n+1},\frac1n\right)
  \ \cup\ (1,+\infty)
\]
è un'unione di intervalli aperti. L'esempio mostra che l'unione infinita
dei singoletti chiusi $\{1/n\}$ non è necessariamente chiusa, mentre
l'aggiunta del suo punto di accumulazione $0$ la rende tale.

\begin{schema}
\begin{tikzpicture}[x=10cm,y=1cm]
  \foreach \quota/\nome in {.6/X,-.6/Y} {
    \draw[asse] (-.045,\quota) -- (1.065,\quota);
    \node[anchor=east] at (-.075,\quota) {$\nome$};
    \foreach \n in {1,...,9} {
      \fill[coloreschema] ({1/\n},\quota) circle[radius=1.7pt];
    }
    \node[coloreschema,inner sep=0pt] at (.047,\quota) {$\cdots$};
  }
  \draw[estremo] (0,.6) circle[radius=2.5pt];
  \fill[coloreschema] (0,-.6) circle[radius=2.5pt];
  \foreach \pos/\testo in {0/0,.25/{\frac14},.333333/{\frac13},.5/{\frac12},1/1} {
    \node[above=6pt] at (\pos,.6) {$\testo$};
  }
  \node[below=5pt] at (0,-.6) {$0$};
  \node[below=5pt] at (1,-.6) {$1$};
\end{tikzpicture}
\par\smallskip
I punti $1/n$ si avvicinano a $0$: il cerchio vuoto indica che
$0\notin X$, il punto pieno che $0\in Y$. Aggiungendo $0$ si ottiene
l'insieme chiuso $Y=X\cup\{0\}$.
\end{schema}

\section{Intorni di un punto}
% Fonte: Lezione 8, pp.~10--11

\paragraph{Definizione generale.}
Dato $x_0\in\R$, un insieme $U\subseteq\R$ si dice \emph{intorno di $x_0$}
se esiste un aperto $A$ tale che
\[
  x_0\in A\subseteq U.
\]
L'insieme $U$ non deve necessariamente essere aperto: deve soltanto
contenere un insieme aperto che comprenda $x_0$. Per esempio,
$[x_0-1,x_0+1]$ è un intorno di $x_0$, perché contiene
$(x_0-1,x_0+1)$.

Se $U$ è esso stesso aperto e $x_0\in U$, si parla di \emph{intorno aperto}
di $x_0$. In particolare, ogni intervallo $(a,b)$ con $a<x_0<b$
è un intorno aperto di $x_0$. L'insieme di tutti gli intorni di $x_0$
si può indicare con $\mathcal I_{x_0}$.

La definizione precedente equivale a una forma molto concreta:
\[
  U\text{ è un intorno di }x_0
  \quad\Longleftrightarrow\quad
  \exists\,\varepsilon>0:\ (x_0-\varepsilon,x_0+\varepsilon)\subseteq U.
\]
Infatti, se $x_0\in A\subseteq U$ con $A$ aperto, il criterio di apertura
fornisce un tale $\varepsilon$; nel verso opposto, l'intervallo stesso
può fare da insieme aperto $A$.

\paragraph{Intorni sferici.}
Per $r>0$ si definisce l'\emph{intorno sferico}
di centro $x_0$ e raggio $r$ come
\[
  I_r(x_0)=(x_0-r,x_0+r)
  =\{x\in\R:|x-x_0|<r\}.
\]
Sulla retta si tratta semplicemente di un intervallo aperto simmetrico.
Ogni intorno $U$ di $x_0$ contiene un intorno sferico $I_r(x_0)$.
Per esempio, se $(a,b)\subseteq U$ e $a<x_0<b$, basta scegliere
$0<r\le\min\{x_0-a,b-x_0\}$.

\begin{schema}
\begin{tikzpicture}[x=1.2cm,y=1cm]
  \draw[asse] (-.5,0) -- (8.5,0);
  \draw[draw=black!25,line width=4pt] (0,0) -- (8,0);
  \draw[draw=black!55,fill=white,line width=.9pt] (0,0) circle[radius=2.6pt];
  \draw[draw=black!55,fill=white,line width=.9pt] (8,0) circle[radius=2.6pt];
  \draw[intervallo] (1.4,0) -- (4.6,0);
  \foreach \pos in {1.4,4.6} {
    \draw[estremo] (\pos,0) circle[radius=2.6pt];
    \draw[densely dotted,black!45] (\pos,.1) -- (\pos,.65);
  }
  \fill[coloreschema] (3,0) circle[radius=2pt];
  \draw[densely dotted,black!45] (3,.1) -- (3,.65);
  \draw[{Stealth[length=1.5mm]}-{Stealth[length=1.5mm]},coloreschema]
    (1.4,.55) -- node[above] {$r$} (3,.55);
  \draw[{Stealth[length=1.5mm]}-{Stealth[length=1.5mm]},coloreschema]
    (3,.55) -- node[above] {$r$} (4.6,.55);
  \foreach \pos/\testo in {0/a,1.4/{x_0-r},3/{x_0},4.6/{x_0+r},8/b} {
    \node[below=6pt] at (\pos,0) {$\testo$};
  }
  \node[black!70] at (6.5,.5) {$(a,b)\subseteq U$};
\end{tikzpicture}
\par\smallskip
L'intorno sferico $I_r(x_0)$, in blu, è simmetrico rispetto a $x_0$
e resta dentro $(a,b)$; quindi è contenuto anche in $U$.
\end{schema}

\begin{proposizione}[Aperti e intorni]
Un insieme $A\subseteq\R$ è aperto se e solo se è un intorno di ciascuno
dei suoi punti.
\end{proposizione}
\begin{proof}
Se $A$ è aperto e $x\in A$, nella definizione di intorno possiamo
scegliere come aperto proprio $A$. Viceversa, se $A$ è un intorno di
ogni $x\in A$, per ogni tale $x$ esiste $\varepsilon_x>0$ con
$(x-\varepsilon_x,x+\varepsilon_x)\subseteq A$; questo è esattamente
il criterio che definisce un insieme aperto.
\end{proof}

\paragraph{Intorni di $+\infty$ e $-\infty$.}
Quando si lavora nella retta estesa
$\overline\R=\R\cup\{-\infty,+\infty\}$, gli intorni aperti di
$+\infty$ contengono una coda $(a,+\infty]$, e quelli di $-\infty$
contengono una coda $[-\infty,b)$, con $a,b\in\R$.
Se si resta nella sola $\R$, le rispettive code si scrivono
$(a,+\infty)$ e $(-\infty,b)$: i simboli $\pm\infty$ non sono allora
punti dell'insieme ambiente.

\section{Intorni laterali}
% Fonte: Lezione 8, p.~12

Fissato $x_0\in\R$, un \emph{intorno sinistro}
di $x_0$ è un intervallo $(a,x_0)$ con $a<x_0$; un \emph{intorno destro}
è un intervallo $(x_0,b)$ con $b>x_0$. Si possono ammettere anche
$a=-\infty$ e $b=+\infty$.\footnote{Gli intorni laterali qui definiti
escludono il punto $x_0$: non sono quindi intorni di $x_0$ nel senso
della definizione generale. La denominazione serve a descrivere il
comportamento vicino al punto da sinistra o da destra.}

Per esempio, $(-1,1)$ è un intorno sferico di $0$ di raggio $1$.
È anche un intorno ordinario di ogni suo punto, fra cui $0$, $1/2$,
$1/3$, $-1/10$ e $1/5$. Visto invece rispetto ai suoi estremi, $(-1,1)$ è un
intorno destro di $-1$ e un intorno sinistro di $1$. In particolare,
gli intorni laterali di un punto non contengono quel punto.

\begin{schema}
\begin{tikzpicture}[x=1.15cm,y=1cm]
  \foreach \quota in {.7,-.7} {
    \draw[asse] (0,\quota) -- (8,\quota);
    \draw[densely dotted,black!40] (4,{\quota-.25}) -- (4,{\quota+.3});
    \node[below=6pt] at (4,\quota) {$x_0$};
  }
  \draw[intervallo] (1,.7) -- (4,.7);
  \draw[intervallo] (4,-.7) -- (7,-.7);
  \foreach \pos/\quota in {1/.7,4/.7,4/-.7,7/-.7} {
    \draw[estremo] (\pos,\quota) circle[radius=2.6pt];
  }
  \node[below=6pt] at (1,.7) {$a$};
  \node[below=6pt] at (7,-.7) {$b$};
  \node[above=6pt,coloreschema] at (2.5,.7) {intorno sinistro $(a,x_0)$};
  \node[above=6pt,coloreschema] at (5.5,-.7) {intorno destro $(x_0,b)$};
\end{tikzpicture}
\par\smallskip
Negli intorni laterali il punto $x_0$ è escluso,
come mostrano i cerchi vuoti.
\end{schema}
